\section{Expected Results and Broader Impacts}
\label{sec:results_impacts}

\subsection{Expected Intellectual and Technical Outcomes}
The execution of this multi--scale computational framework is projected to yield a series of deterministic, high--resolution transport metrics and open--access tools, replacing empirical approximations with first--principles physics. The specific technical deliverables include:
\begin{itemize}
    \item \textbf{Mechanistic Particulate Trapping Maps:} High--resolution, multi--dimensional profiles of particle--capture efficiency mapping the explicit boundary layer dynamics ($Re \ll 1$) across varied benthic architectures, isolating the exact roles of coral canopy geometries and sand--matrix pore topologies under shear.
    \item \textbf{A Priori Transport Tensors:} A comprehensive database of mathematically derived, effective hydrodynamic dispersion tensors ($\bm{D}_{\text{eff}}$) and continuous kinetic trapping rates ($R$) calculated across a broad parameter space of localized salinity, temperature, and shear profiles characteristic of the Tayrona National Natural Park (TNNP) bays.
    \item \textbf{Validated Open--Source Solvers:} A modular, high--performance computational suite integrating discrete Stokesian Dynamics with continuum Advection--Diffusion--Reaction (ADR) mapping frameworks, made publicly available via GitHub to enable reproducible multi--scale source--code auditing by the global fluid mechanics community.
    \item \textbf{Peer--Reviewed Academic Outputs:} The dissemination of foundational transport findings through high--impact physical and computational engineering journals, directly bridging the gap between theoretical microstructural rheology and applied marine conservation.
\end{itemize}

\subsection{Broader Impacts and Regional Territorial Value}
Beyond its contributions to fundamental mathematical physics, this research establishes a vital technical bridge to the ecological and socio--economic management of the Caribbean Colombian coastline:
\begin{itemize}
    \item \textbf{Grounding Regional Environmental Forensics:} By providing mathematically rigorous parameters to replace \emph{ad hoc} constants, this framework directly enhances the predictive fidelity of macro--scale coastal models (e.g., ROMS, Delft3D) utilized by local institutes such as INVEMAR. This allows environmental engineers to accurately forecast the structural degradation of reef buffers and microplastic dispersion patterns under shifting industrial climate stressors.
    \item \textbf{Fostering Regional Computational Capacity:} The deployment of advanced multi--scale numerical frameworks within the Magdalena department serves as an educational catalyst. By introducing engineering students and regional researchers to the integration of low--level platform execution, discrete particle mechanics, and high--performance computing (UPR--HPCf), this project actively addresses the curricular gap between abstract transport theory and numerical high--performance implementation.
%   \item \textbf{Strengthening the Caribbean Tech Ecosystem:} This project actively builds public technical capacity within regional developer networks (such as local open--source foundations). By demonstrating that high--performance scientific computing can be leveraged as an open--source territorial asset, the research fosters a highly specialized, localized community of systems engineers capable of addressing complex regional environmental challenges from a rigorous, handmade computational baseline.
%   \item \textbf{Rectifying Macro--Scale Model Vulnerabilities under Shifting Climate Regimes:} Current regional--scale hydrodynamic and environmental forecasting frameworks (e.g., ROMS, Delft3D) are fundamentally constrained by their spatial grid resolution, which compresses complex benthic boundary layers into flat, uniform interfaces. To compensate, these macro--models rely heavily on empirical, parameter--fitted drag and mass--transfer coefficients. Because these classic tuning parameters are inherently incapable of adapting to non--linear physical changes, this project provides the critical, first--principles missing link. By deriving the explicit mathematical dependencies of $\bm{D}_{\text{eff}}$ and $R$ directly from multi--body microstructure mechanics, this work provides the exact upscaled physical parameters required to stabilize macro--scale predictions, transforming regional environmental forecasting from an empirical fitting exercise into a rigorous, predictive science under unprecedented climate stress.
    \item \textbf{Establishing a Methodological Paradigm Shift Toward True Multi--Scale Coastal Modeling:} As identified in the foundational limitations of current large--scale architectures, macrostructural forecasting suites (e.g., ROMS, Delft3D) operate under an inherent closure problem, relying on static, empirical boundary parameters that fail under non--linear climate disruptions. This project explicitly acknowledges that a universal adjustment of global coastal models requires accounting for a wide array of distinct boundary interfaces (e.g., mangrove networks, macrostructural cliffs, varying sediment types). Rather than attempting a brute--force, universal parameterization, this study serves as the critical, first--principles \emph{demonstration case} and methodological blueprint. By rigorously isolating and upscaling the complex, living benthic canopy of the Tayrona region using discrete Stokesian Dynamics, this work establishes the reproducible computational and mathematical pipeline required to close the macrostructural transport equations. It marks a definitive first step away from empirical data--fitting models and toward a mathematically rigorous, predictive multi--scale environmental science.
\end{itemize}
